% Source: https://artofproblemsolving.com/community/c3h2085046_quadratic__completing_the_square_versus_quadratic_equation
% Stable identity: aops-topic-2085046-q01
Completing the square: \(x^2+2x=\dfrac{1}{2}\), so \((x+1)^2=\dfrac{3}{2}\) and \(x=-1\pm\sqrt{\dfrac{3}{2}}\).

Quadratic formula: \(x=\dfrac{-4\pm\sqrt{16+8}}{4}=\dfrac{-4\pm 2\sqrt6}{4}=\dfrac{-2\pm\sqrt6}{2}\).

They agree, because rationalising gives \(\sqrt{\dfrac{3}{2}}=\dfrac{\sqrt3}{\sqrt2}=\dfrac{\sqrt6}{2}\), so \(-1\pm\sqrt{\dfrac{3}{2}}=\dfrac{-2\pm\sqrt6}{2}\). Both methods are correct:
\[
\boxed{x=\dfrac{-2\pm\sqrt6}{2}}
\]
